\begin{center}
	{\fontsize{24pt}{24pt}\selectfont\textbf{\textcolor{black}{Lời cảm ơn}}}
\end{center}

Để hoàn thành đồ án này, trước hết em xin gửi lời cảm ơn sâu sắc đến tập thể phòng Lab 209B3 và thầy Bùi Quốc Bảo, người đã trực tiếp hướng dẫn em trong luận văn tốt nghiệp này, cũng như thầy, cô giảng viên trong Khoa Điện – Điện tử, Trường Đại học Bách Khoa – ĐH Quốc gia TP. HCM, 
đã truyền đạt kiến thức và kinh nghiệm quý báu cho em trong suốt quá trình học tập và rèn luyện của em tại trường. 
Trong quá trình làm luận văn, em đã gặp không ít khó khăn, còn thiếu sót nhiều về mặt kiến thức thực tế. Nhưng với sự giúp đỡ, hỗ trợ từ anh Cao Văn Hiếu, anh Huỳnh Thanh Sang trong việc thiết kế phần cứng và những người bạn tại phòng Lab 209B3 trong việc chia sẻ kiến thức với em, và đặc biệt là sự giúp đỡ,
tư vấn nhiệt tình từ thầy Bùi Quốc Bảo chia sẻ những kiến thức và kinh nghiệm của thầy để em có thể hoàn thành tốt luận văn của mình và tích lũy thêm cho bản thân được nhiều những kinh nghiệm thực tế, kiến thức hữu ích cho bản thân, 
cải thiện kỹ năng tìm kiếm thông tin trong quá trình thực tập để hoàn thành nhiệm vụ được thầy giao. Dù bản thân đã rất cố gắng để hoàn thành tốt các nhiệm vụ được giao nhưng không thể tránh khỏi những sai sót nên em rất mong nhân được sự thông cảm, đóng góp ý kiến của quý thầy cô.\\

Cuối cùng, em xin kính chúc quý thầy cô sức khỏe, luôn thành công trong công việc và cuộc sống. Em xin chân thành cảm ơn!\\
\begin{flushright}
    \begin{tabular}{@{}c@{}}
        TP.Hồ Chí Minh, ngày 15 tháng 5 năm 2026 \\
        \\[1.5cm] % Chừa khoảng trống 1.5cm để ký tên
        \textbf{Nguyễn Xuân Phát} % \textbf để in đậm tên
    \end{tabular}
\end{flushright}